% !TeX document-id = {ac384373-0975-4c4e-b6b2-65611dc82067}
% !TeX TXS-program:compile = txs:///pdflatex/[--shell-escape]
\documentclass[a4paper,12pt]{article}

\usepackage[T2A]{fontenc}
\usepackage[russian]{babel}
\usepackage[intlimits,namelimits]{amsmath}
\usepackage{mathtools}

\usepackage{indentfirst}
\usepackage{icomma}
\usepackage{graphicx}
\usepackage{geometry}
\usepackage{setspace}
\usepackage{fancyhdr}
\usepackage{titlesec}

\usepackage{minted}

\usepackage[bibstyle=gost-numeric,citestyle=numeric-comp]{biblatex}

\usepackage{hyperref}

\addbibresource{biblio.bib}

\geometry{a4paper, nomarginpar, left=20mm, right=10mm, top=20mm, bottom=20mm}

\setlocalecaption{russian}{ref}{Список использованной литературы}

\newcommand\Vect[1]{{\boldsymbol{#1}}}
\newcommand{\dd}{\mathrm{d}}
\newcommand{\diag}{\operatorname{diag}}
\renewcommand{\Re}{\operatorname{Re}}
\renewcommand{\Im}{\operatorname{Im}}
\renewcommand{\exp}{\operatorname{e}}
\renewcommand{\imath}{\mathrm{i}}

\newcommand{\mintedlabel}[1]{\phantomsection\label{\detokenize{#1}}}

% \DeclarePairedDelimeter\bra{\langle}{|}
\DeclarePairedDelimiter\ket{|}{\rangle}
\DeclarePairedDelimiter\abs{\lvert}{\rvert}
\DeclarePairedDelimiterX\commut[2]{[}{]}{#1,#2}

\titlespacing{\subsection}{\parindent}{3.5ex plus 1ex minus .2ex}{2.3ex plus .2ex}
\titleformat{\section}{\normalfont\Large\bfseries}{\thesection.}{1em}{}
\titlespacing{\section}{\parindent}{3.5ex plus 1ex minus .2ex}{2.3ex plus .2ex}

\fancypagestyle{plain}{%
  \fancyhf{}%
  \fancyfoot[R]{\thepage}%
  \renewcommand{\headrulewidth}{0pt}%
  \renewcommand{\footrulewidth}{0pt}%
}

%%%%%%%%%%%%%%%%%%%%%%%%%%%%%%%%%%%%%%%%%%%%%%%%%%%%%%%%%%%%%%%%%%%%%%%%%%%%%%%%

\makeatletter
\def\@sect[#1]#2{%
	\refstepcounter{app@section}
	\setcounter{equation}{0}
	\addcontentsline{toc}{section}{Приложение~\theapp@section}
	{%
		{
			\raggedleft%
			\normalfont%
			Приложение~\theapp@section\par%
		}%
		\Large\bfseries%
		#2%
		\par%
	}
	\vskip 3ex
}
% This is supposed to be "starred" version, but it doesn't have sense in the
% appendix, so we simply repeat above definition.
\def\@ssect#1{%
	\refstepcounter{app@section}
	\addcontentsline{toc}{section}{Приложение~\theapp@section}%
	{%
		{%
			\raggedleft%
			\normalfont%
			Приложение~\theapp@section\par%
		}%
		\Large\bfseries%
		#1%
		\par%
	}%
	\vskip 3ex
}

\renewcommand\appendix{%
	\newpage%
	\newcounter{app@section}
	\setcounter{app@section}{0}
	%\setcounter{equation}{0}%
	\renewcommand\theapp@section{\Asbuk{app@section}}
	\renewcommand\theequation{\theapp@section\arabic{equation}}
	%%% We need to redefine \section command here...
	%%% NB: this is messy, VERY messy (see GOST for autoreferat!)
	\renewcommand{\section}{%
		\secdef\@sect\@ssect
	}
}

%%%%%%%%%%%%%%%%%%%%%%%%%%%%%%%%%%%%%%%%%%%%%%%%%%%%%%%%%%%%%%%%%%%%%%%%%%%%%%%%

\begin{document}
\fontsize{14pt}{20pt}\selectfont	

\onehalfspacing
\pagestyle{plain}

\begin{titlepage}
  \begin{center}
    {%
      \scshape
      Министерство науки и высшего образования Российской Федерации\\
    }
    федеральное государственное бюджетное образовательное учреждение\\
    высшего образования \\
    «Иркутский государственный университет»\\
    (ФГБОУ ВО «ИГУ»)
  \end{center}
  \vspace*{3em}

  \begin{flushright}
    \begin{minipage}{70mm}
      Физический факультет\\
      Кафедра теоретической физики\\
      Зав. кафедрой, к.ф.-м.н., доцент\\
      \underline{\hspace{20mm}} Ловцов С.В.
    \end{minipage}
  \end{flushright}
  \vspace*{4em}

	\begin{center}
    КУРСОВАЯ РАБОТА\\[2ex]\
    \bfseries%
		ОЦЕНКА КАЧЕСТВА ЧИСЛЕННОГО РЕШЕНИЯ УРАВНЕНИЯ ОСЦИЛЛЯЦИЙ НЕЙТРИНО В 
		СРЕДЕ
	\end{center}

  \vspace*{3em}

  \begin{flushright}
    \begin{minipage}{80mm} 
    Руководитель: к.ф.-м.н., доцент\\
    \underline{\hspace{30mm}}Ломов В.П.\\
    Студент гр. 01211-ДБ\\
    \underline{\hspace{20mm}}/Кучеров Даниил Андреевич\\
    Работа защищена с\\
    оценкой \underline{\hspace{40mm}}\\
    <<\underline{\hspace{8mm}}>>\underline{\hspace{22mm}}2024г.\\
    Нормконтролёр:
    \underline{\hspace{30mm}}/Фролова А.С./\\
    \end{minipage}
  \end{flushright}

  \vfill
  \begin{center}
    Иркутск - 2024
  \end{center}
\end{titlepage}

\setcounter{page}{2}

\thispagestyle{empty}
\tableofcontents

\newpage

\section*{Введение}\addcontentsline{toc}{section}
{Введение}

Нейтрино (итал. neutrino — нейтрончик, уменьшительное от neutrone — нейтрон) —
общее название нейтральных фундаментальных частиц с полуцелым спином,
участвующих только в слабом и гравитационном взаимодействиях, и относящихся к
классу лептонов.

Нейтрино необычная частица, так как представлена как бы в двойном виде:
массивное и флейворное состояния. Массивное нейтрино отвечает за распространение
в пространстве. В свою же очередь, флейворное нейтрино отвечает за
взаимодействие. В общем случае нейтринные поля, обладающие определенным
флейвором, являются суперпозицией полей нейтрино обладающей определенной
массой. Флейворные состояния нейтрино не идентичны массовым состояниям. Данная
суперпозиция осуществляется с помощью унитарной матрицы смешивания \(U\):
\begin{equation*}
  \ket{\nu_{\alpha}}=\sum_{i}U_{\alpha i}^{*}\ket{\nu_{i}},
\end{equation*}
где \(\ket{\nu_{i}}\) — вектор состояния поля массивных нейтрино, соответственно
\(\ket{\nu_{\alpha}}\) — вектор состояния поля флейворных нейтрино.

Также справедливо и обратное утверждение, что поля массивных нейтрино являются
суперпозицией полей флейворных нейтрино, оно осуществляется с помощью
эрмитого сопряженной матрицы:
\begin{equation*}
  \ket{\nu_{i}}=\sum_{\alpha}U_{\alpha i}\ket{\nu_{\alpha}}
\end{equation*}

Главной причиной для построения такой необычной картины мира послужил тот факт,
что в экспериментах по регистрации солнечных нейтрино их получалось меньше, чем
предсказывала теория. Для объяснения этого факта пришлось потребовать наличия
массы у нейтрино и ввести два типа нейтрино: флейворное и массивное. Таким
образом получилось разрешить парадокс недостатка солнечных нейтрино, а само
явление получило название \textit{нейтринных осцилляций}.

Исследования в области нейтринных осцилляций являются одним из главных
источников информации о нейтрино. В настоящее время работает большое число
нейтринных экспериментов, которые детектируют события от приходящих нейтрино,
проходящих по мере своего распространения как через вакуум, так и через
вещество, что необходимо учитывать при анализе собранных данных.

В данной работе мы проведем анализ уравнения нейтринных осцилляций, а также сравним два метода интегрирования уравнения осцилляций нейтрино в среде. Первый метод использует широко известную схему Рунге-Кутта, второй --- предназначен специально для матричных дифференциальных уравнений.   

\newpage
\section{Уравнение нейтринных осцилляций в среде}

Когда активные флэйворные нейтрино распространяются в веществе, на их
эволюционное уравнение влияют эффективные потенциалы из-за когерентного
взаимодействия со средой посредством реакций нейтрального и заряженного
токов. Таким образом, влияние среды можно описать в терминах потенциалов, в
которых распространяются нейтрино, зависящим от состава среды, электрической
нейтральности, намагниченности (ориентации спинов), скоростей частиц среды.

Зная, что состояния флейворных нейтрино являются суперпозицией состояний
массовых нейтрино
\begin{equation}\label{eq:0}
	\ket{\nu_{\alpha}}=\sum_{i}U_{\alpha i}^{*}\ket{\nu_{i}},
\end{equation} 
матрица \(U_{\alpha i}\) для нейтрино Дирака выражается в виде
\begin{equation}
	U=O_{23}\Gamma O_{13}\Gamma^{\dagger}O_{12},
\end{equation}
где \(O_{ij}\) ортогональные матрицы, представляющие повороты в пределах
\([0,\pi/2]\) в соответствующих плоскостях,
\(\Gamma=\diag(1,1,\exp^{\imath\delta})\)---диагональная матрица, содержащая
фазу нарушения CP-инвариантности \(\delta\), она меняется в пределах
\([0,2\pi]\).

Рассмотрим нейтрино \(\nu_{\alpha}\) возникшее в среде в момент времени
\(t_{0}\). Состояние системы \(\ket{\Psi(t)}\) в момент времени \(t\geq t_{0}\)
можно выразить как
\begin{equation}
	\ket{\Psi(t)}=\sum_{\beta}\Psi_{\beta}(t)\ket{\nu_{\beta}},
\end{equation}
где \(\Psi_{\beta}(t_{0})=\delta_{\alpha\beta}\). Вероятность нахождения
флейвора \(\beta\) в точке \(r=t\) (Здесь и ниже мы используем единицы
\(\hbar=c=1\)) равна
\begin{equation}\label{eq:8}
	P_{\alpha\beta}(r)=|\Psi_{\beta}(r)|^{2}.
\end{equation}
При выходе нейтрино из среды амплитуды начинают изменятся в соответствии с
уравнением колебаний в вакууме, решение которого становится гораздо проще, если
записать его через массовые состояния. Используя уравнение \eqref{eq:0}, а также
обозначив через \(A_{j}=\phi_{j}(r_{*})\) амплитуду вероятности нахождения
\(\ket{\nu_{j}}\) на краю среды, при \(r\geq r_{*}\) мы можем записать
\begin{equation}\label{eq:7}
	\Psi_{\beta}(r)=\sum_{j=1}^{3}U_{\beta j}A_{j}\exp^{-\imath E_{j}L}
\end{equation}
где \(E_{j}=\sqrt{|\Vect{p}|^{2}+m_{j}^{2}}\), \(L=r-r_{*}\) это расстояние
пройденное нейтрино в вакууме. Теперь мы подставляем выражение \eqref{eq:7} в
уравнение \eqref{eq:8} мы получим
\begin{equation}
	P_{\alpha\beta}=\sum_{j}|U_{\beta j}|^{2}|A_{j}|^{2}+2\sum_{i>j}\Re[U_{\beta i}U_{\beta j}^{*}A_{i}A_{j}^{*}\exp^{-\imath\Delta_{ij}L}].
\end{equation}
Величина \(\Delta_{ij}=\Delta m_{ij}/2E\), где \(E=|\Vect{p}|\), представляет
собой волновое число колебаний связанное с квадратичной разностью масс
\(\Delta m_{ij}^{2}=m_{i}^{2}-m_{j}^{2}\).

В результате задача вычисления \(P_{\alpha\beta}\) сводится к нахождению
амплитуд собственных массовых состояний \(\phi_{j}(r)\) внутри среды (а именно
\(A_{j}\)), при начальном условии \(\phi_{j}(r_{0})=U_{\alpha j}^{*}\).

Для релятивийских нейтрино, распространяющихся в среде, после вычитания
глобальной фазы, эволюционное уравнение для амплитуд имеет вид
\begin{equation}\label{eq:1}
	\imath\frac{\dd\Phi}{\dd\xi}=[H_{0}+U^{\dagger}VU]\Phi,
\end{equation}
где \(\Phi^{T}(r)=(\phi_{1}(\xi),\phi_{2}(\xi),\phi_{3}(\xi))\) — это амплитуда
вероятностей массовых состояний нейтрино, \(\xi=\frac{r}{R_{0}}\), где \(R_{0}\) --- солнечный радиус, \(H_{0}\) обозначает гамильтониан для
эволюции в вакууме, он имеет вид:
\begin{equation}
	H_{0}=\frac{a}{E}
  \begin{pmatrix}
    0 & 0 & 0\\
    0 & b & 0\\
    0 & 0 & 1
  \end{pmatrix},
  \quad
  a=4,35196\times10^{6},
  \quad
  b=0,030554,
\end{equation}
\(E\) — численное значение энергии нейтрино выраженное в MeV, \(U\) — матрица
смешивания, \(V=\diag(V_{cc},0,0)\), где \(V_{cc}\) — потенциал среды, зависящий от расстояния, который можно выразить как
\begin{equation}\label{eq:12}
	\nu(r)=\sqrt{2}G_{F}n_{e}(r),
\end{equation}
\(G_{F}\) --- постоянная Ферми, \(n_{e}(r)\) --- плотность электронов на пути нейтрино. Из выражения \eqref{eq:12} можно вывести функции эффективных плотностей (мы следуем работе \cite{Casas:2016asi})
\begin{equation}\label{eq:10}
	\nu(\xi)=\gamma\exp^{-\eta\xi}
	\quad
	\gamma=6,5956\times10^{4},
	\quad
	\eta=10,54
	\quad 
	\xi=\frac{r}{R_{0}} ,
\end{equation}
\begin{equation}\label{eq:11}
	\nu(\xi)=\frac{\gamma}{\xi^{3}}
	\quad
	\gamma=52,934
\end{equation}
для солнечной и для сверхновой моделей соответственно.

Уравнение \eqref{eq:1} можно переписать в более удобной форме. Если использовать преобразование \(\Psi(\xi)=\Gamma^{\dag}\Phi(\xi)\) (\(\Gamma\) была введена
выше).  Затем, вынести \(V_{cc}(\xi)=\nu(\xi)\)  и ввести обозначение для матрицы \(W=U^{\dag}\diag(1,0,0)U\), то мы получаем
уравнение:
\begin{equation}\label{eq:9}
  \imath\frac{\dd\Psi}{\dd\xi}=[H_{0}+\nu(\xi)W]\Psi,
\end{equation}
Матрица \(W\) имеет вид:
\begin{equation}
  W=
  \begin{pmatrix}
    c_{13}^{2}c_{12}^{2} & c_{12}s_{12}c_{13}^{2} & c_{12}c_{13}s_{13}\\
    c_{12}s_{12}c_{13}^{2} & s_{12}^{2}c_{13}^{2} & s_{12}c_{13}s_{13}\\
    c_{12}s_{13}c_{13} & s_{12}c_{13}s_{13} & s_{13}^{2}
  \end{pmatrix}.
\end{equation}
\(c_{12}\), \(c_{13}\) — это значения косинусов углов смешивания, а \(s_{12}\),
\(s_{13}\) — это значение синусов углов смешивания, их значения были взяты из
источника \cite{Casas:2016asi}.

Уравнение \eqref{eq:1} является матричным, комплекснозначным, причем
\(H(\xi)=[H_{0}+\nu(\xi)W]\) это эрмитова матрица гамильтониана
эволюции. Решением является волновая функция \(\Psi\). Затем, мы расчитываем
вероятность выживания частицы по формуле
\begin{equation}
	\langle P_{ee} \rangle=c_{12}^{2}c_{13}^{2}|\Psi_{1}|^{2}+s_{12}^{2}c_{13}^{2}|\Psi_{2}|^{2}+s_{13}^{2}|\Psi_{3}|^{2}.
\end{equation}
На данный момент уравнение имеет только численное решение, которое требует
массивных численных интегрирований линейной однородной системы обыкновенных
дифференциальных уравнений (ОДУ) с коэффициентами, зависящими от расстояния,
которое нейтрино проходят вдоль среды.

\newpage
\section{Метод численного интегрирования}
Для численного интегрирования систем дифференциальных уравнений могут использоваться различные методы. Достаточно распространенным из них является метод схем Рунге-Кутта, в частности схемы 4-5 и 5-6 порядков точности. Главным недостатком метода Рунге-Кутта является рассмотрение каждого уравнения системы отдельно от остальных, что при интегрировании уравнения \eqref{eq:9} приводит к нарушению нормировки решения.

Существуют и другие методы численного интегрирования, в частности, методы геометрического интегрирования, ниже рассмотрим один из таких.

Рассмотрим оператор эволюции \(U(t,t_{0})\), который переводит
квантовомеханическое состояние (волновую функцию) \(\Psi\) от времени \(t_{0}\)
ко времени \(t\)
\begin{equation*}
	\Psi(t)=U(t,t_{0})\Psi(t_{0}),
\end{equation*}
откуда следует очевидное условие \(U(t_{0},t_{0})=I\), где \(I\) — единичный
оператор.

С течением времени волновая функция \(\Psi\) не должна менять свою норму — это
константное значение, обеспечивающее сохранение вероятности. Математически это
означает, что оператор \(U(t,t_{0})\) является унитарным
\begin{equation*}
	U^{\dag}(t,t_{0})U(t,t_{0})=U(t,t_{0})U^{\dag}(t,t_{0})=I
\end{equation*}
и удовлетворяет уравнению (см. также \cite{Shaydurova:2019})
\begin{equation}\label{eq:2}
	\imath\hbar\frac{\dd U(t,t_{0})}{\dd t}=\lambda H(t)U(t,t_{0}),
\end{equation}
где \(H(t)\) — гамильтониан системы, который в общем случае зависит от времени,
а \(\lambda\) — промежуточный параметр, который в конце полагается равным
единице.  Если бы мы не рассматривали матричное уравнение, то решение можно было
бы найти с помощью итераций
\begin{equation}\label{eq:3}
	U(t,t_{0})=I+\sum_{n=1}^{\infty}\lambda^{n}P_{n}(t,t_{0}),\qquad P_{n}(t,t_{0})=0
\end{equation}

Подставляя выражение \eqref{eq:3} в уравнение \eqref{eq:2}, и приравнивая
слагаемые при равных степенях \(\lambda\), возникает система уравнений на
\(P_{n}\)
\begin{equation*}
		\imath\hbar\frac{\dd P_{1}(t,t_{0})}{\dd t}= H(t),\qquad
			\imath\hbar\frac{\dd P_{n}(t,t_{0})}{\dd t}= H(t)P_{n-1}(t,t_{0})\quad
			(n>1)
\end{equation*}
интегрируем, получаем \(P_{n}\)
\begin{equation}\label{eq:4}
	P_{n}(t,t_{0})=\Big(-\frac{\imath}{\hbar}\Big)^{n}
  \int_{t_{0}}^{t}dt_{1}\int_{t_{0}}^{t_{1}}dt_{2}\ldots
  \int_{t_{0}}^{t_{n-1}}dt_{n}H_{1}H_{2}\ldots H_{n},\qquad
	H_{i}=H(t_{i}).
\end{equation}

Если гамильтониан не зависит от времени, то \(P_{n}\) принимает достаточно
компактный вид и ряд \eqref{eq:3} сводится к экспоненте
\begin{equation}
	P_{n}=\frac{1}{n!}\bigg(-\imath(t-t_{0})\frac{H}{\hbar}\bigg)^{n}\qquad
	\Longrightarrow\quad
	U(t,t_{0})=\exp^{-\imath(t-t_{0})\lambda\frac{H}{\hbar}}
\end{equation}

Но в нашем случае мы ищем решение матричного уравнения, для того чтобы выражение
\eqref{eq:4} было верным, необходимо чтобы \(\commut{H(t_{1})}{H(t_{2})}=0\),
\(\forall t_{1},t_{2}.\) В общем случае это не так для \(H=H(t)\). По этому мы и
не можем решить аналитически уравнение нейтринных осцилляций.

Наконец, рассмотрим метод численного геометрического интегрирования, а именно метод 
разложения Магнуса. 
Этот метод относится к так
называемым геометрическим интеграторам. Отличительной особенностью разложения
Магнуса является то, что приближенное решение, предоставляемое процедурой,
разделяет с точным решением соответствующие геометрические свойства. Более
конкретно, для рассматриваемого случая соответствующее свойство соответствует
унитарному характеру оператора эволюции времени, следовательно, норме вектора
решения. В численном расчёте это проверяется на каждом шаге интегрирования, и,
следовательно, вероятность сохраняется за счёт окончательного решения,
полученного путём составления последовательных шагов. Сам метод представляет
собой систематический способ построения аппроксимаций решения нестационарного
уравнения Шрёдингера таким образом, чтобы в любом порядке оператор эволюции был
унитарным.

Пусть мы имеем дифференциальное уравнение вида:
\begin{equation*}
	\frac{dY(t)}{dt}=A(t)Y(t)\qquad
	A\in C^{N\times N},\quad
	Y\in C^{N}.
\end{equation*}
Решение этого уравнения используя метод Магнуса будет иметь вид:
\begin{equation*}
	Y(t)=\exp^{\Omega(t,t_{0})}Y_{0},\qquad
	\Omega(t,t_{0})\in C^{N\times N},\quad
	\Omega(t_{0},t_{0})=0,
\end{equation*}
где \(\Omega(t,t_{0})\) — это ряд Магнуса:
\begin{equation*}
	\Omega=\sum_{k=1}^{\infty}\lambda^{k}\Omega_{k}.
\end{equation*}
Первые три члена ряда Магнуса имеют вид:
\begin{equation*}
  \begin{lgathered}
    \Omega_{1}(t,t_{0})=\int_{t_{0}}^{t} dt_{1}A_{1},\\
    \Omega_{2}(t,t_{0})=\frac{1}{2}\int_{t_{0}}^{t} dt_{1}
    \int_{t_{0}}^{t_{1}} dt_{2}\commut{A_{1}}{A_{2}},\\
    \Omega_{3}(t,t_{0})=\frac{1}{6}\int_{t_{0}}^{t} dt_{1} \int_{t_{0}}^{t_{1}}
    dt_{2} \int_{t_{0}}^{t_{2}} dt_{3}\Big(
    \commut[\big]{A_{1}}{\commut{A_{2}}{A_{3}}}+
    \commut[\big]{\commut{A_{1}}{A_{2}}}{A_{3}}\Big).
  \end{lgathered}
\end{equation*}
Для дальнейших вычислений нам будет достаточно первых двух слагаемых. 

\newpage
\section{Оценка качества интегрирования уравнения}

Рассмотрим два способа для интегрирования уравнения \eqref{eq:9} : метод схем Рунге-Кутта и метод разложения Магнуса --- и сравним их.
Далее, вычисления будут проводится для солнечной модели, энергию нейтрино берем равную 1 MeV, все исходные значения
взяты из статьи \cite{Casas:2016asi}

Критерием оценивания качества решений, полученных методами, изложенными выше, будет являться условие сохранения нормировки решения, то есть
\begin{equation}
  \label{eq:5}
  \sum_{i}^{3}|\Psi_{i}|^{2}-1=0.
\end{equation}

Для реализации первого способа используем стандартную процедуру из Wolfram Mathematica (мы пытаемся повторить результаты работы изложенной в \cite{Casas:2016asi}, однако в статье приведены не все необходимые параметры интегрирования).
Второй способ реализуем, используя алгоритм М4 приведенный в статье \cite{Casas:2016asi}. Анализ полученных значений осуществляем в python. Ниже приведены графики
функции контроля нормировки (рис.\ref{fig:rk-wolf} и рис.\ref{fig:magnus})
\begin{equation}
  \label{eq:6}
  F(\xi)=\sum_{i}^{3}|\Psi_{i}(\xi)|^{2}-1.
\end{equation}

\begin{figure}[htb]
  \centering%
  \includegraphics[width=0.8\linewidth]{Math}
  \caption{\label{fig:rk-wolf}График контроля нормировки для решения полученного
    с помощью Wolfram Mathematica}
\end{figure}

\begin{figure}[htb]
  \centering%
  \includegraphics[width=0.8\linewidth]{magexp}
  \caption{\label{fig:magnus}График контроля нормировки для решения полученного
    с помощью метода разложения Магнуса}
\end{figure}

На графиках видно, что в случае Mathematica функция в начале имеет тенденцию к
линейному убыванию, а затем — значительный рост, причем с ненулевой первой
производной, достигая максимального значения магнитуды в области интегрирования
\(0,05\). В случае метода Магнуса, мы видим близкое к линейному убывание, вот
только порядок значений магнитуды \(10^{-10}\), что гораздо меньше чем в
Mathematica. Отсюда можно сделать вывод, что метод интегрирования по Магнусу
гораздо лучше сохраняет нормировку решения, чем Mathematica.

\clearpage
\section*{Заключение}\addcontentsline{toc}{section}
{Заключение}

В работе было оценено качество интегрирования уравнения осцилляций нейтрино в
среде с помощью метода Магнуса и метода Рунге-Кутта.  Как следует из
представленных данных, метод разложения Магнуса намного лучше справляется с
задачей интегрирования уравнения нейтринных осцилляций, чем метод
Рунге"--~Кутта. На то есть несколько причин: во-первых, метод Магнуса, в отличие
от Рунге"--~Кутта, он учитывает свойства уравнения в целом, гарантируя
сохранение нормировки решения, а значит и точность интегрирования. Вторая
замечательная особенность метода Магнуса состоит в динамическом шаге, с помощью
которого можно гораздо лучше описывать быстро осциллирующие функции.

Подводя итог, можно сказать, что методы Рунге"--~Кутта демонстрируют малую
эффективность для интегрирования уравнения нейтринных осцилляций по сравнению с
методом интегрирования Магнуса.

\clearpage

\addcontentsline{toc}{section}{Список использованной литературы}
\printbibliography

\newpage

\appendix

\section{Численное интегрирование в Mathematica}
\label{app:sec:worlfram}

В Mathematica мы численно интегрируем уравнение \eqref{eq:9}. Необходимые
параметры инициализируем в начале программы.
\begin{minted}[linenos,fontsize=\footnotesize,escapeinside=||]{nb}
  Y = 6.5956*10^4
  X = 10.54
  v[x_] := Y Exp[-X x]
  t = 1
  a = 4.35196*10^6
  b = 0.03055
  H0 = a/t {
    {0, 0, 0},
    {0, b, 0},
    {0, 0, 1}
  }
  s12 = Sqrt[0.308]
  s13 = Sqrt[0.0234]
  c12 = Sqrt[1 - s12^2]
  c13 = Sqrt[1 - s13^2]
  W = {
    {c13^2*c12^2, c12*s12*c13^2, c12*c13*s13},
    {c12*s12*c13^2, s12^2*c13^2, s12*c13*s13},
    {c12*s13*c13, s12*c13*s13, s13^2}
  }
  f[x_] := {f1[x], f2[x], f3[x]}
  
  s3 = NDSolve[{                                                             |\label{code:ndsolve11}|
    I*f'[x] == (H0 + v[x] W) . f[x],
    f1[0.1] == c12*c13, f2[0.1] == s12*c13, f3[0.1] == s13},
    {f1, f2, f3}, {x, 0.1, 0.5}]                                             |\label{code:ndsolve12}|

  psi4[y_] = Evaluate[{f1[x], f2[x], f3[x]} /. s3] /. {x -> y} /. {w_} -> w  |\label{code:eval1}|

  s4 = NDSolve[{                                                             |\label{code:ndsolve21}|
    I*f'[x] == (H0 + v[x] W) . f[x],
    f1[0.1] == c12*c13, f2[0.1] == s12*c13, f3[0.1] == s13},
    {f1, f2, f3}, {x, 0.1, 1}]                                               |\label{code:ndsolve22}|

  psi5[y_] = Evaluate[{f1[x], f2[x], f3[x]} /. s4] /. {x -> y} /. {w_} -> w  |\label{code:eval2}|

  tab12 = Table[{                                                            |\label{code:tab11}|
    x,
    Re[psi4[x][[1]]], Im[psi4[x][[1]]],
    Re[psi4[x][[2]]], Im[psi4[x][[2]]],
    Re[psi4[x][[3]]], Im[psi4[x][[3]]],
    Arg[psi4[x][[1]]],   Arg[psi4[x][[2]]],   Arg[psi4[x][[3]]],
    Abs[psi4[x][[1]]]^2, Abs[psi4[x][[2]]]^2, Abs[psi4[x][[3]]]^2,
    Abs[psi4[x][[1]]]^2 + Abs[psi4[x][[2]]]^2 + Abs[psi4[x][[3]]]^2 - 1,
    c12^2 c13^2 Abs[psi4[x][[1]]]^2 +
    s12^2 c13^2 Abs[psi4[x][[2]]]^2 +
    s13^2 Abs[psi4[x][[3]]]^2},
    {x, 0.101, 0.5, 0.001}];                                                 |\label{code:tab12}|

  Export["Mathpsi4.dat", tab12, "Table"]                                     |\label{code:export1}|

  tab13 = Table[{                                                            |\label{code:tab21}|
    x,
    Re[psi5[x][[1]]], Im[psi5[x][[1]]],
    Re[psi5[x][[2]]], Im[psi5[x][[2]]],
    Re[psi5[x][[3]]], Im[psi5[x][[3]]],
    Arg[psi5[x][[1]]],   Arg[psi5[x][[2]]],   Arg[psi5[x][[3]]],
    Abs[psi5[x][[1]]]^2, Abs[psi5[x][[2]]]^2, Abs[psi5[x][[3]]]^2,
    Abs[psi5[x][[1]]]^2 + Abs[psi5[x][[2]]]^2 + Abs[psi5[x][[3]]]^2 - 1,
    c12^2 c13^2 Abs[psi5[x][[1]]]^2 +
    s12^2 c13^2 Abs[psi5[x][[2]]]^2 +
    s13^2 Abs[psi5[x][[3]]]^2},
    {x, 0.101, 1, 0.001}];                                                   |\label{code:tab22}|

  Export["Mathpsi5.dat", tab13, "Table"]                                     |\label{code:export2}|
\end{minted}

В строках \ref{code:ndsolve11}--\ref{code:ndsolve12} и
\ref{code:ndsolve21}--\ref{code:ndsolve22} используем стандартную функцию
численного интрегрирования \mintinline{nb}{NDSolve} с автоматическим выбором
параметров (расчет проводился в Mathematica 13.2.1.0).

Для промежуточных проверок мы проводили расчет на половине интервала и на всем
интервале \((0,1{,}\,1)\).

Для сравнения численных расчетов результаты численного интегрирования
сохраняются в файл (строки \ref{code:export1} и \ref{code:export2}).

Из-за особенности вывода функции \mintinline{nb}{NDSolve} нам приходится сначала
конвертировать результат в функциональную форму (строки \ref{code:eval1} и
\ref{code:eval2}), а затем вычислять по ним численные данные (строки
\ref{code:tab11}--\ref{code:tab12} и \ref{code:tab21}--\ref{code:tab22}).

Странный формат строк \ref{code:eval1}, \ref{code:eval2} связан с процессом
подстановки (код с \mintinline{nb}{/.}) и формальным определением функции
(переменная \mintinline{nb}{y_}).

Нужно заметить, что результат работы \mintinline{nb}{NDSolve} — это
\mintinline{nb}{InterpolatingFunction} — интерполяционный полином, поэтому
остается неизвестным, как результаты расчёта связаны с действительным численным
интегрированием, поскольку в строках \ref{code:tab11}--\ref{code:tab12},
\ref{code:tab21}--\ref{code:tab22} мы фактически перевычисляем значения в точках
заданных точках, в которых, возможно, не производилось численное интегрирование
уравнения.

\newpage
\section{Численное интегрирование по методу Магнуса}

Для численного интегрирования по методу Магнуса использовался
код\footnote{Использовалась версия с номером коммита
  \texttt{d3f0634b63a848aa36b4854464ae083c0f97dfcf}.} из проекта magexp
(\url{https://git.sr.ht/~vp1981/magexp}). Программа
запускалась как
\begin{minted}[fontsize=\footnotesize]{bash}
  (\
    export LC_NUMERIC=en_US.UTF-8; \
    for b in $(seq 0.101 1e-3 1); \
    do \
      ./m4-tol -c "b=${b}" >> data_10.dat;\
    done\
  )
  sed -i '/^#/d' data_10.dat
\end{minted}
с конфигурацией
\begin{minted}[linenos,fontsize=\footnotesize]{lua}
  tol = 1e-3
  model = "sun"
  E = 1e0
  a = 0.1
  b = 1.0
  s12 = math.sqrt(0.308)
  c12 = math.sqrt(1 - s12*s12)
  s13 = math.sqrt(0.0234)
  c13 = math.sqrt(1 - s13*s13)
  psi0 = { {c12*c13,0}, {s12*c13,0}, {s13,0} }
\end{minted}

\newpage
\section{Анализ и визуализация данных}
\label{app:sec:python}

Для анализа и визуализации результатов численного интегрирования уравнения
осцилляций мы использовали библиотеки numpy и matplotlib.

Ниже приведен код вспомагательных функций, использованных для виртуализации.%
\inputminted[linenos,fontsize=\footnotesize,escapeinside=||]{python}{draw-funcs.py}%
Сам анализ производился в интерактивной сессии jupyter (графическая сессия
qtconsole).

Графики выводились в отдельном окне (инструкция \texttt{\%matplotlib qt}),
вспомагательный код подключался через инструкцию \texttt{\%load}.

Из-за особенностей записи данных (по колонкам) и количеству данных, требовалась
специальная загрузка данных
\begin{minted}[breaklines,fontsize=\footnotesize]{python}
p2,      r12, i12, r22, i22, r32, i32, ph12, ph22, ph32, a12, a22, a32, c2, pr2     = np.loadtxt("Mathpsi5.dat", unpack=True)
p1, pr1, r11, i11, r21, i21, r31, i31, ph11, ph21, ph31, a11, a21, a31, c11, c12, _ = np.loadtxt("data_10.dat", unpack=True)
\end{minted}
Последняя цифра в имени переменной указывает на источник данных, имя — на тип
данных.

Примеры вызова функций для построения графиков:
\begin{minted}[fontsize=\footnotesize]{python}
draw1b(p1, c21, labx=r'$\xi$', laby=labs['c1'], title='Контроль magexp')
draw1b(p2,  c2, labx=r'$\xi$', laby=labs['c1'], title='Контроль Mathematica')
\end{minted}
\end{document}

%%% Local Variables:
%%% mode: LaTeX
%%% TeX-master: t
%%% TeX-command-extra-options: "-shell-escape"
%%% End:
